\documentclass[11pt]{article}
\usepackage[margin=1in]{geometry}
\usepackage{amsmath,amssymb,amsthm,mathtools}
\usepackage{array}
\usepackage[hidelinks]{hyperref}

\newtheorem{theorem}{Theorem}[section]
\newtheorem{lemma}[theorem]{Lemma}
\newtheorem{proposition}[theorem]{Proposition}
\newcommand{\spE}{s^+}
\newcommand{\smE}{s^-}
\newcommand{\sig}{\sigma}
\newcommand{\one}{\mathbf 1}

\title{Positive Square Energy of Every Pentacyclic Cactus}
\author{}
\date{26 July 2026}

\begin{document}
\maketitle

\begin{abstract}
Let $G$ be a connected pentacyclic cactus of order $n$.  We prove
$s^+(G)>n$.  The exact block-additive doubly nonnegative estimate reduces the
problem to the cycle multisets $\{3,3,3,3,q\}$ and $\{3,3,3,5,5\}$.
For a disconnected shared-cut graph, we give the complete lists of eleven and
fifteen colored cluster partitions.  The two partitions not settled by direct
packet addition require entry-sensitive interval decompositions of a central
pentagon; the all-singleton partitions require reduced-tree leaf arguments.
For one shared-cut cluster, the first residual is proved by opening the
$q$-cycle or partitioning it among four triangular petals.  The second is
settled by an exact census of all forty colored incidence trees: thirty-six
admit a positive one-cycle split, and four explicit exceptions are treated
directly.  Every construction allows arbitrary connector trees and arbitrary
trees attached at arbitrary vertices.
\end{abstract}

\section{Preliminaries and dependencies}

All graphs are finite, simple, and undirected.  A cactus is a connected graph
whose blocks are bridges or cycles.  Its cyclomatic number is the number of
cyclic blocks.  If the adjacency eigenvalues of $G$ are
$\lambda_1,\ldots,\lambda_n$, put
\[
 \spE(G)=\sum_{\lambda_j>0}\lambda_j^2,
 \qquad \smE(G)=\sum_{\lambda_j<0}\lambda_j^2,
 \qquad \sig(G)=\spE(G)-|V(G)|.
\]
For every vertex partition $V(G)=V_1\mathbin{\dot\cup}\cdots
\mathbin{\dot\cup}V_k$, induced-subgraph superadditivity gives
\begin{equation}\label{eq:super}
 \spE(G)\geq\sum_{i=1}^k\spE(G[V_i]).
\end{equation}
This is the theorem of Akbari, Kumar, Mohar, and Pragada~\cite{AKMP}.

Write $T=C_3$, $P=C_5$, and
\[
 \delta=\sec(\pi/5)-1=\sqrt5-2<\frac12,
 \qquad \delta_q=\sec(\pi/q)-1<1\quad(q\equiv1\pmod4).
\]
We use the following established packet bounds.  All permit arbitrary trees
attached at arbitrary vertices; mixed packets may have arbitrary bridge
connectors.
\begin{align}
 \sig(T)&>0,& \sig(P)&\geq-\delta,
   &\sig(C_q)&\geq-\delta_q,\label{eq:packet-one}\\
 \sig(TT)&>1,& \sig(TP)&>1-\delta,
   &\sig(TC_q)&>1-\delta_q,\label{eq:packet-two}\\
 \sig(H)&\geq0&&\text{for every bicyclic or tricyclic cactus }H,
   \label{eq:packet-small}\\
 \sig(TTT)&>2&&\text{for one shared-cut triangular cluster},
   \label{eq:packet-ttt}\\
 \sig(TTTT)&>3&&\text{for one shared-cut four-triangle cluster},
   \label{eq:packet-tttt}\\
 \sig(TTC_q)&>2-\delta_q
   &&\text{if the two triangles share a cut and }q\equiv1\pmod4,
   \label{eq:packet-ttq}\\
 \sig(TPP)&>6-2\sqrt5
   &&\text{for one shared-cut cluster}.\label{eq:packet-tpp}
\end{align}
Also, every connected tetracyclic cactus has positive surplus:
\begin{equation}\label{eq:tetracyclic}
 \sig(H)>0.
\end{equation}
Bounds \eqref{eq:packet-one}--\eqref{eq:packet-small} are proved in the
bicyclic and tricyclic companion manuscripts~\cite{Bicyclic,Tricyclic}.
The shared $TTC_q$ phase comparison \eqref{eq:packet-ttq} is the
intersecting-triangle case of the $\{3,3,q\}$ theorem in~\cite{Tricyclic}.
If every cycle is $3\pmod4$ and the cycle-packing number is at most two, the
Sachs half-plane theorem~\cite{PackingTwo} gives
\begin{equation}\label{eq:favorable}
 \sig(H)>r-1
\end{equation}
for a connected $r$-cyclic cactus.  This proves
\eqref{eq:packet-ttt}.  For \eqref{eq:packet-tttt}, packing at most two again
uses \eqref{eq:favorable}; the sole packing-three shared-cut incidence is a
central triangle with three pairwise disjoint triangular petals.  Its direct
Sachs calculation proves $\spE>\smE$, hence
$\spE>|E|=|V|+3$.  This calculation, including the exact matching injections
which dominate the three-cycle Sachs term, appears in
\cite[Proposition~4.11]{PackingTwo}.  Bound \eqref{eq:packet-tpp} is the exact
shared $C_3,C_5,C_5$ phase bound in~\cite{Tricyclic};
\eqref{eq:tetracyclic} is the tetracyclic cactus theorem~\cite{Tetracyclic}.
These are the only external packet dependencies used below.

\section{Cluster trees and exact territories}

Join two cyclic blocks in the \emph{shared-cut graph} when they share a cut
vertex, and call its connected components shared-cut clusters.  In the
block-cut tree of $G$, contract the incidence subtree belonging to each
shared-cut cluster to one marked node.  Every cluster, including every
singleton cyclic block, is marked.  Take the minimal subtree spanning all
marked nodes and suppress only unmarked degree-two nodes.  The resulting tree
is the \emph{reduced cluster tree}.  A reduced edge therefore represents a
path in the block-cut tree whose block nodes are all bridge blocks; it contains
no unassigned cyclic block.

\begin{lemma}[Arbitrary connector territories]\label{lem:territories}
Let $S_1,\ldots,S_k$ be pairwise vertex-disjoint connected subtrees of the
reduced cluster tree, and suppose that every marked cluster node belongs to
exactly one $S_i$.  Then $V(G)$ has a partition into connected induced
territories $G_1,\ldots,G_k$ such that $G_i$ retains exactly the cyclic blocks
whose marked nodes lie in $S_i$.  All cuts between territories occur on
actual bridge edges.  Every tree attached anywhere to the cyclic hull is
assigned whole to one territory.
\end{lemma}

\begin{proof}
First contract each $S_i$ to a labelled root in the reduced cluster tree.
Assign every remaining reduced-tree vertex to a nearest labelled root, using
a fixed label order to break ties.  In a tree, the vertices assigned to one
root induce a connected subtree: along the unique path toward that root the
distance decreases by one, while no competing distance can become strictly
smaller after the fixed tie rule.  Thus this assignment partitions the
reduced tree into connected labelled subtrees.  Every edge between two labels
is a reduced edge and hence expands to a path consisting only of bridge blocks.
Choose one actual bridge on that path as the cut.  Assign the two path remnants
to their respective labels.

Equivalently, at an unmarked Steiner branch one may assign the branch vertex
and one initial segment to one incident territory and cut one actual bridge on
every other incident segment.  This explicitly includes a genuine $Y$-shaped
connector and does not assume a linear order of clusters.  After all selected
bridge cuts, each component of the cyclic hull is connected and contains
exactly its assigned cyclic blocks.

Finally, a component outside the cyclic hull meets it at a unique vertex.  Two
attachments would create a cycle in the block-cut tree.  Assign the entire
hanging component to the territory containing that attachment.  The resulting
vertex sets induce connected subgraphs and retain precisely the asserted
cycles.  Edges crossing between vertex sets are discarded only through
\eqref{eq:super}; no edge-monotonicity statement is used.
\end{proof}

We will also split a cyclic block into proper consecutive intervals.  If the
shared vertices of a cycle $C$ are listed cyclically and the node $C$ is
deleted from its cycle--cut incidence tree, the resulting branches can be
assigned nonempty consecutive intervals of $C$, one interval containing the
shared vertex of each branch.  Every interval induces a proper path, so the
split cycle contributes no cycle to any territory.  Several cyclically
adjacent marks may instead be put in one interval.  Hanging trees are assigned
by their unique attachment as in Lemma~\ref{lem:territories}.

\section{The exact DNN residual reduction}

For completeness, we record the exact reduction rather than merely citing its
output.  For a connected edge-containing graph define
\[
 \kappa(G)=\sup_{\substack{M\succeq0,\ M\geq0\\
                         \one^TM\one>0}}
 \frac{4\bigl(\sum_{uv\in E(G)}\sqrt{M_{uv}}\bigr)^2}
      {\one^TM\one}.
\]
The sharp cactus DNN theorem~\cite{DNN,LTZ} states that, if $G$ has $b$
bridge blocks and cyclic blocks of lengths $\ell_1,\ldots,\ell_r$, then
\begin{equation}\label{eq:kappa}
 \kappa(G)=b+\sum_{j=1}^r\kappa(C_{\ell_j}),\qquad
 \kappa(C_\ell)=
 \begin{cases}
  \ell,&\ell\text{ even},\\
  \ell\sec^2(\pi/(2\ell)),&\ell\text{ odd},
 \end{cases}
 \qquad \smE(G)\leq\kappa(G).
\end{equation}
Here is the proof mechanism behind the last inequality and the exact
additivity.  Fold every nonedge entry $M_{uv}$ into the two diagonal entries
by adding $M_{uv}(e_u-e_v)(e_u-e_v)^T$.  This preserves the denominator and
the edge entries.  Rotational averaging on $C_\ell$, concavity of the square
root, and the least Fourier eigenvalue give the two cycle constants in
\eqref{eq:kappa}; equality is attained by
$x(2cI+A(C_\ell))$, where $c=1$ for even $\ell$ and
$c=\cos(\pi/\ell)$ for odd $\ell$.

For a one-vertex sum $G=G_1\vee_zG_2$, use a folded Gram system and decompose
$g_z=u+w+h$, where $u$ and $w$ lie in the mutually orthogonal spans of the two
sides and $h$ is orthogonal to both.  Use $u+h$ as the cut vector on the first
side and $w$ on the second.  Both the edge numerator sum and the denominator
split; Cauchy--Schwarz gives
$\kappa(G)\leq\kappa(G_1)+\kappa(G_2)$.  Conversely, glue near-optimal Gram
systems orthogonally and scale them so that
$T_1/\kappa(G_1)=T_2/\kappa(G_2)$, proving equality.  Finally write the
adjacency matrix as $A=A_+-B$.  The Schur product theorem makes $B\circ B$
doubly nonnegative, while
\[
 \smE/2=-\sum_{uv\in E(G)}B_{uv}
 \leq\sum_{uv\in E(G)}|B_{uv}|,
 \qquad \one^T(B\circ B)\one=\operatorname{tr}B^2=\smE.
\]
The definition of $\kappa$ yields $(\smE)^2\leq\kappa(G)\smE$ and hence the
last assertion of \eqref{eq:kappa}.

Put
\[
 \epsilon_\ell=\begin{cases}
 0,&\ell\text{ even},\\
 \ell\tan^2(\pi/(2\ell)),&\ell\text{ odd}.
 \end{cases}
\]
Then
\begin{equation}\label{eq:epsilon}
 \epsilon_3=1,
 \qquad \epsilon_5=5-2\sqrt5,
 \qquad \epsilon_5+\epsilon_7<1,
\end{equation}
and $\epsilon_\ell$ strictly decreases through odd lengths.  The monotonicity
follows, after putting $u=\pi/(2x)$, from the positivity of the derivative of
$\tan^2u/u$, since $2u>\sin u\cos u$.

For a pentacyclic cactus with cycle lengths $\ell_1,\ldots,\ell_5$, block
counting gives
\[
 b+\sum_{i=1}^5\ell_i=n+4.
\]
Since $\spE+\smE=2|E|=2(n+4)$, \eqref{eq:kappa} gives the exact residual
estimate
\begin{equation}\label{eq:dnn}
 \sig(G)=\spE(G)-n\geq4-\sum_{i=1}^5\epsilon_{\ell_i}.
\end{equation}
If at most two cycles are triangles, then
$\sum_i\epsilon_{\ell_i}\leq2+3\epsilon_5<4$.  With exactly three
triangles, \eqref{eq:epsilon} shows that the sum can reach $4$ only when both
remaining cycles are pentagons: replacing either pentagon by an even cycle or
an odd cycle of length at least seven makes their combined contribution less
than one.  With at least four triangles, designate four of them and call the
fifth cycle $Q$; this includes $Q=T$.  Consequently \eqref{eq:dnn} proves
$\sig(G)>0$ unless the cycle multiset belongs to exactly one of the residual
families
\begin{equation}\label{eq:residual}
 TTTTQ=\{3,3,3,3,q\}\quad(q\geq3),
 \qquad TTTPP=\{3,3,3,5,5\}.
\end{equation}

\section{Disconnected shared-cut graph: the $TTTTQ$ family}

\begin{proposition}\label{prop:disconnected-ttttq}
Suppose the five cyclic blocks are $TTTTQ$ and the shared-cut graph is
disconnected.  Then $\spE(G)>|V(G)|$.
\end{proposition}

\begin{proof}
If $q\equiv1\pmod4$, the complete list of proper colored set partitions,
up to permutation of the four triangles, consists of the following eleven
rows.  Each row records an induced-territory certificate.  The cluster
territories are realized by Lemma~\ref{lem:territories}; a displayed zero is
nonnegative, while a strict triangular term makes the total strict.
\begin{center}
\renewcommand{\arraystretch}{1.14}
\begin{tabular}{>{\ttfamily}l l}
partition & lower surplus certificate\\ \hline
TTTT|Q & $>3-\delta_q>0$\\
TTTQ|T & positive tetracyclic territory plus a strict $T$\\
TTT|TQ & $>2+(1-\delta_q)>0$\\
TTQ|TT & nonnegative tricyclic territory plus $>1$\\
TTT|T|Q & $>2+0-\delta_q>0$\\
TT|TT|Q & $>2-\delta_q>0$\\
TTQ|T|T & nonnegative tricyclic territory plus two strict $T$'s\\
TT|TQ|T & $>1+(1-\delta_q)+0$\\
TT|T|T|Q & $>1+0+0-\delta_q>0$\\
TQ|T|T|T & $>1-\delta_q$ plus strict $T$'s\\
T|T|T|T|Q & reduced-tree leaf argument below
\end{tabular}
\end{center}
The table uses \eqref{eq:packet-one}--\eqref{eq:tetracyclic}; in particular
the connected shared-cut $TTT$ and $TTTT$ clusters satisfy
\eqref{eq:packet-ttt} and \eqref{eq:packet-tttt}.

For the final row, inspect the minimal reduced tree spanning the five marked
cycle nodes.  Every leaf of this minimal tree is marked.  If a triangle is a
leaf, cut the first actual bridge toward the rest.  Its territory has strict
positive surplus, and the connected remainder is tetracyclic and has strict
positive surplus by \eqref{eq:tetracyclic}.  If no triangle were a leaf, every
leaf would have to be the unique $Q$ node, impossible because a finite
nontrivial tree has at least two leaves.  This proves the all-singleton row
without assuming that the reduced tree is a path.

If $q$ is even, an even unicyclic territory has zero surplus.  Apply the same
cluster territories: every territory with at most three cycles is
nonnegative, every tetracyclic territory is positive, and the proper cluster
partition always leaves either a strict triangular territory or a positive
tetracyclic territory.  The all-singleton leaf proof is unchanged.  If
$q\equiv3\pmod4$, every cluster with at most three cycles has packing number
at most two and is strict by \eqref{eq:favorable}; every four-cycle territory
is strict by the tetracyclic theorem \eqref{eq:tetracyclic}.  Singleton
territories are strict as well.  Thus these two nonhostile parities are also
covered.
\end{proof}

\section{Disconnected shared-cut graph: the $TTTPP$ family}

The direct cluster ledger leaves two entry-sensitive partitions.  We prove
them before displaying the complete audit.

\begin{lemma}[The $TTTP|P$ repair]\label{lem:tttp-p}
Suppose one shared-cut cluster $A$ has cycles $T_1,T_2,T_3,P_0$, and a remote
pentagon $P_1$ is bridge-separated from it.  Then $\spE(G)>|V(G)|$.
\end{lemma}

\begin{proof}
If two triangles of $A$ meet, let $k$ be the number of shared cyclic cuts on
$P_0$.  The incidence excess for four cycle nodes is three.  The $k$ cuts on
$P_0$ contribute at least $k$, and the intersecting triangle pair contributes
one additional unit, whether its common cut is off $P_0$ or is a cut of degree
at least three on $P_0$.  Hence $k\leq2$.  Choose a private vertex $v$ of
$P_0$, put $v$ and all off-cycle branches rooted there in a tree $F$, and put
the rest of $A$ in $H$.  Then $P_0-v$ connects all of its shared cuts, so $H$
is connected.  Its only cycles are the three triangles, including an
intersecting pair; hence their packing number is at most two.  Thus
\[
 \sig(A)\geq\sig(F)+\sig(H)>-1+2=1.
\]
Together with $\sig(P_1)\geq-\delta$, connector territories give total
surplus $>1-\delta>0$.

Assume now that the three triangles are pairwise disjoint.  Incidence-tree
acyclicity forces $P_0$ to be central, with $T_i$ attached at three distinct
vertices $x_i$.  Let the connector from $P_1$ first meet the cyclic core of
$A$ at $z$.

If it enters through $T_i$, including at $x_i$, put all of $T_i$, the
connector, and $P_1$ in one territory.  Put $P_0-x_i$ and the other two
triangles in the other.  These territories have cycle multisets $TP$ and
$TT$.  An entry through a branch rooted on $T_i$ is the same case.

If the connector enters $P_0$ at a private vertex $z$, cyclically order
$z,x_1,x_2,x_3$.  Choose an $x_i$ adjacent to $z$ in this cyclic order and
cut the two boundary edges of a consecutive interval containing exactly the
pair $z,x_i$ among these four marks.  This interval, $T_i$, the connector,
and $P_1$ form a $TP$ territory.  The complementary proper interval and the
other two triangles form a $TT$ territory.  All branches rooted at a boundary
vertex go with the territory owning that vertex.

In every entry topology, the surplus is
\[
 >(1-\delta)+1>0.
\]
The split of $P_0$ is essential: it prevents the common attachment vertex from
being assigned to two packets.
\end{proof}

\begin{lemma}[The $TTP|T|P$ repair]\label{lem:ttp-t-p}
Suppose one shared-cut cluster $A$ has cycles $T_1,T_2,P_0$, while the
singleton clusters are $B=T_3$ and $C=P_1$.  Then
$\spE(G)>|V(G)|$.
\end{lemma}

\begin{proof}
If $T_1,T_2$ share a cut, \eqref{eq:packet-ttq} gives
$\sig(A)>2-\delta$; with
$\sig(T_3)>0$ and $\sig(P_1)\geq-\delta$, the total is
$>2-2\delta>0$.  Thus assume that the incidence in $A$ is the distinct-cut
chain
\[
 T_1-P_0-T_2,
\]
where $T_i$ meets $P_0$ at $x_i$.

Use the actual minimal connector tree on the three marked clusters.  If the
$B$--$C$ path avoids $A$, take the entire $B$--$C$ connector subtree as a
$TP$ territory and $A$ as the other territory.  Their surplus is
$>1-\delta+0$.  This includes a genuine Steiner $Y$: assign its branch vertex
to the $B$--$C$ territory and cut the $A$ branch on an actual bridge.

Suppose instead that the $B$--$C$ path passes through $A$, so there are two
external connector entries into its cyclic core.  An entry through $T_i$,
including at $x_i$, is represented by a mark at $x_i$ forced to travel with
$T_i$; an entry at a private vertex of $P_0$ is marked at that vertex.

If both entries are forced to the same $T_i$, put both external cycles, both
connector branches, and all of $T_i$ in one $TTP$ territory.  Put $P_0-x_i$
and the other triangle in a strict triangular territory.  This also covers
coincident entries at $x_i$.

Suppose the external pentagon entry is forced to $T_i$, but the external
triangle entry is not.  Isolate $x_i$ by two cuts of $P_0$.  Its interval
carries $P_1+T_i$ and is a $TP$ territory; the complementary interval carries
$T_3+T_j$ and is a $TT$ territory.  If instead the external triangle entry is
forced to $T_i$, isolate $x_i$ on the $TT$ side and put $P_1+T_j$ on the
complementary $TP$ side.

It remains that both entries lie at private vertices of $P_0$.  If they
coincide, take a proper interval containing that common mark and one adjacent
$x_i$.  It gives a $TTP$ territory, while the complement gives a triangular
territory.  If the entries are distinct, cyclically order the two labelled
external marks and the two unlabeled marks $x_1,x_2$.  One neighbor of the
pentagon-entry mark is an $x_i$ such that the other $x_j$ is adjacent to the
triangle-entry mark; this follows by checking the three cyclic orders of two
labelled and two unlabeled marks.  Cut in the two gaps between these adjacent
pairs.  The resulting proper intervals carry $P_1+T_i$ and $T_3+T_j$, hence
give $TP$ and $TT$ territories.

Every case therefore gives either
\[
 TP+TT,\qquad \sig>(1-\delta)+1>0,
\]
or
\[
 TTP+T,\qquad \sig>0,
\]
where the tricyclic packet is nonnegative and the triangle is strict.  This
covers private cycle entries, shared-cut entries, entries through attached
triangles, coincident roots, and branching connectors.
\end{proof}

\begin{proposition}\label{prop:disconnected-tttpp}
Suppose the five cyclic blocks are $TTTPP$ and the shared-cut graph is
disconnected.  Then $\spE(G)>|V(G)|$.
\end{proposition}

\begin{proof}
Up to permutation of equal cycles, the complete list of fifteen proper colored
set partitions is the following.  The first and seventh rows invoke
Lemmas~\ref{lem:tttp-p} and \ref{lem:ttp-t-p}; all other non-singleton rows use
the direct packet ledger.
\begin{center}
\renewcommand{\arraystretch}{1.14}
\begin{tabular}{>{\ttfamily}l l}
partition & certificate or lower surplus\\ \hline
TTTP|P & Lemma~\ref{lem:tttp-p}\\
TTT|PP & $>2+0$\\
TTPP|T & positive tetracyclic territory plus a strict $T$\\
TTP|TP & $>0+(1-\delta)$\\
TPP|TT & $>(6-2\sqrt5)+1$\\
TTT|P|P & $>2-2\delta$\\
TTP|T|P & Lemma~\ref{lem:ttp-t-p}\\
TT|TP|P & $>1+(1-\delta)-\delta$\\
TT|T|PP & $>1+0+0$\\
TPP|T|T & $>6-2\sqrt5$ plus strict $T$'s\\
TP|TP|T & $>2(1-\delta)$ plus a strict $T$\\
TT|T|P|P & $>1-2\delta$ plus a strict $T$\\
TP|T|T|P & $>1-2\delta$ plus strict $T$'s\\
PP|T|T|T & nonnegative $PP$ plus strict $T$'s\\
T|T|T|P|P & reduced-tree leaf argument below
\end{tabular}
\end{center}
Every displayed quantity is positive because $\delta<1/2$.

For the all-singleton row, inspect the minimal reduced tree spanning the five
marked cycle nodes.  Every leaf is marked.  If a triangle is a leaf, cut its
first actual bridge; the triangular territory is strict, and the connected
tetracyclic remainder is strict by \eqref{eq:tetracyclic}.  If no triangle is
a leaf, all leaves are among the two pentagons.  There must be at least two
leaves, so both pentagons are leaves and there are exactly two leaves.  A finite
tree with exactly two leaves is a path, with all three triangle marks internal.
Cut actual bridges so that the path is partitioned into $TP$, $T$, and $TP$
territories.  Their total surplus is
$>2(1-\delta)>0$.  This argument remains valid when reduced edges expand to
long paths or pass through previously suppressed degree-two vertices.
\end{proof}

\section{Fully shared $TTTTQ$}

For one shared-cut cluster, let $I$ be the bipartite incidence tree whose cycle
nodes are the five cyclic blocks and whose cut nodes are the vertices belonging
to at least two cyclic blocks.  If there are $c$ cut nodes, then
\begin{equation}\label{eq:incidence-excess}
 |E(I)|=c+4,
 \qquad \sum_{x\text{ cut node}}(\deg_I(x)-1)=4.
\end{equation}
A multiway cut incident with $d$ cyclic blocks contributes exactly $d-1$.

\begin{proposition}[Fully shared $TTTTQ$]\label{prop:shared-ttttq}
Suppose the five blocks $T_1,T_2,T_3,T_4,Q=C_q$, $q\geq3$, form one
shared-cut cluster.  Then $\spE(G)>|V(G)|$.
\end{proposition}

\begin{proof}
First suppose the four designated triangles are pairwise vertex-disjoint.
Every cut node incident with a triangle must also be incident with $Q$, since
it has degree at least two and cannot be incident with another triangle.  A
triangle cannot meet $Q$ at two cuts, because the two triangle--cut--$Q$ routes
would form a cycle in $I$.  Connectedness forces each triangle to meet $Q$
exactly once, and the four attachment cuts are distinct.  Thus $Q$ is central
and the four triangles are pairwise disjoint petals.  This forces $q\geq4$.

List the four attachments in cyclic order and partition $V(Q)$ into four
nonempty consecutive cyclic intervals, one containing each attachment.  Choose
one boundary edge in each arc between consecutive attachments.  Every interval
is a proper path, including when attachments are adjacent; when $q=4$, all
four intervals are singletons.  Adjoin to each interval its triangular petal
and assign every off-core tree by its unique core attachment.  This partitions
$G$ into four connected induced triangular unicyclic cacti, each of strict
positive surplus by \eqref{eq:packet-one}.  Equation \eqref{eq:super} proves
the result.

Now suppose some two designated triangles meet.  Let $k$ be the number of
shared cyclic cuts on $Q$.  Those cuts contribute at least $k$ to
\eqref{eq:incidence-excess}.  The intersecting pair contributes one additional
unit: its common cut is either an additional cut off $Q$, or a cut on $Q$ of
degree at least three and hence contributes at least two.  Therefore
$k+1\leq4$, so $k\leq3$.

Assume $q\geq4$.  Choose a vertex $v$ of $Q$ that is not a shared cyclic cut.
Let $F$ consist of $v$ and every component of $G-E(Q)$ attached to $Q$ only
at $v$, and put $H=G-F$.  Then $F$ is a tree and
\begin{equation}\label{eq:tree-cost}
 \spE(F)=|F|-1.
\end{equation}
The path $Q-v$ contains every shared cut on $Q$; replacing the $Q$ node of $I$
by this path leaves all four triangular blocks connected.  Thus $H$ is
connected and induced, and its only cyclic blocks are the four triangles.

Partition $H$ along bridge connectors between its shared-cut clusters using
Lemma~\ref{lem:territories}.  A cluster containing exactly two intersecting
triangles has surplus $>1$.  One containing three triangles has packing at
most two and surplus $>2$.  A four-triangle cluster has surplus $>3$: use
\eqref{eq:favorable} when its packing number is at most two and
\eqref{eq:packet-tttt} for the central-triangle/three-petal incidence when the
packing number is three.  Every remaining singleton triangular packet has
strict positive surplus.  Hence the cluster containing the chosen intersecting
pair supplies more than one unit and all others supply positive surplus, so
\[
 \spE(H)>|H|+1.
\]
Together with \eqref{eq:tree-cost} and \eqref{eq:super}, this gives
\[
 \spE(G)\geq\spE(H)+\spE(F)
 >|H|+1+|F|-1=|V(G)|.
\]

Finally, if $q=3$, all five cyclic blocks are triangles and the designation of
$Q$ is arbitrary.  The universal triangular block-graph theorem
\cite[Corollary~4.2]{PackingTwo} applies directly and gives
$\spE(G)>|V(G)|$, including all multiway cuts.  The four-petal alternative
cannot have a triangular center because it requires four distinct attachment
vertices.  These alternatives exhaust every incidence.
\end{proof}

\section{Fully shared $TTTPP$: exact incidence census}

Continue to use the incidence tree $I$ and
\eqref{eq:incidence-excess}.  Every cut node has degree at least two, so
$1\leq c\leq4$.  Distinct incidences use distinct vertices of a cycle; hence a
triangle cycle node has degree at most three and a pentagon cycle node degree
at most five.

If an internal cycle node is split into consecutive intervals, deleting it
from $I$ gives one branch multiset for each interval.  Encode such a multiset
by $(t,p)$.  Splitting a triangle leaves total multiset $(2,2)$, and the
following are exactly the positive branch partition types used by the census:
\begin{center}
\renewcommand{\arraystretch}{1.12}
\begin{tabular}{l l}
branches after splitting $T$ & lower surplus\\ \hline
$PP+TT$ & $>1$\\
$T+TPP$ & $>3/2$\\
$TP+TP$ & $>2-2\delta$\\
$P+P+TT$ & $>1-2\delta$\\
$P+T+TP$ & $>1-2\delta$\\
$PP+T+T$ & $>0$
\end{tabular}
\end{center}
Splitting a pentagon leaves total multiset $(3,1)$, with positive types
\begin{center}
\renewcommand{\arraystretch}{1.12}
\begin{tabular}{l l}
branches after splitting $P$ & lower surplus\\ \hline
$P+TTT$ & $>2-\delta$\\
$T+TTP$ & $>0$\\
$TP+TT$ & $>2-\delta$\\
$P+T+TT$ & $>1-\delta$\\
$T+T+TP$ & $>1-\delta$.
\end{tabular}
\end{center}
Here a strict zero means a strict triangular packet is added to a nonnegative
packet.  Every bound follows from
\eqref{eq:packet-one}--\eqref{eq:packet-tpp}, and all are positive because
$\delta<1/2$.

\begin{lemma}[Exact colored incidence census]\label{lem:census}
Modulo permutations of the three triangle nodes, the two pentagon nodes, and
the cut nodes, the numbers of admissible colored incidence trees for
$c=1,2,3,4$ are respectively
\begin{equation}\label{eq:census-counts}
 1,\ 7,\ 18,\ 14.
\end{equation}
Of these, a one-cycle split certified by one of the two preceding tables
resolves respectively
\begin{equation}\label{eq:census-resolved}
 0,\ 6,\ 17,\ 13.
\end{equation}
The four unresolved colored trees are exactly:
\begin{enumerate}
\item one cut $x$ incident with all five cycles;
\item cuts of degrees $(4,2)$, where $x$ meets $T_1,T_2,T_3,P_1$ and $y$
      meets $T_1,P_2$;
\item cuts of degrees $(3,2,2)$, where $x$ meets $T_1,T_2,T_3$, $y$ meets
      $T_1,P_1$, and $z$ meets $T_2,P_2$;
\item a pentagon hub: $P_1$ meets $T_1,T_2,T_3,P_2$ at four distinct
      degree-two cuts.
\end{enumerate}
\end{lemma}

\begin{proof}
The assertion is a finite exact enumeration.  From the repository root it is
reproduced by
\begin{center}
\texttt{python research/}\allowbreak
\texttt{pentacyclic-tttpp-incidence-census.py}.
\end{center}
For each
$c\in\{1,2,3,4\}$ it enumerates subsets of the $5c$ possible cycle--cut edges
having $c+4$ edges.  It rejects a subset unless every cycle degree is positive,
every cut degree is at least two, triangle degrees are at most three, pentagon
degrees are at most five, and union--find verifies the tree condition.

It canonicalizes each surviving edge set over all permutations of the three
triangle labels, both pentagon labels, and all cut labels.  For every cycle
node it then deletes that node, computes the colored branch components by
depth-first search, and compares their sorted $(t,p)$ tuple with the exact
sets
\begin{align*}
 \mathcal G_T={}&\{((0,2),(2,0)),((1,0),(1,2)),((1,1),(1,1)),\\
 &((0,1),(0,1),(2,0)),((0,1),(1,0),(1,1)),
 ((0,2),(1,0),(1,0))\},\\
 \mathcal G_P={}&\{((0,1),(3,0)),((1,0),(2,1)),((1,1),(2,0)),\\
 &((0,1),(1,0),(2,0)),((1,0),(1,0),(1,1))\}.
\end{align*}
These are precisely the two displayed packet tables.  The script asserts the
counts \eqref{eq:census-counts}, the resolved counts
\eqref{eq:census-resolved}, and the following canonical unresolved edge sets,
where cycle nodes $0,1,2$ are triangles, $3,4$ are pentagons, and cut nodes
start at $5$:
\begin{align*}
 c=1:&\quad ((0,5),(1,5),(2,5),(3,5),(4,5));\\
 c=2:&\quad ((0,5),(0,6),(1,5),(2,5),(3,5),(4,6));\\
 c=3:&\quad ((0,5),(0,6),(1,5),(1,7),(2,5),(3,6),(4,7));\\
 c=4:&\quad ((0,5),(1,6),(2,7),(3,5),(3,6),(3,7),(3,8),(4,8)).
\end{align*}
These translate exactly into the four listed incidences.  Thus the census is
both color-preserving and exhaustive.
\end{proof}

\begin{proposition}[Fully shared $TTTPP$]\label{prop:shared-tttpp}
Suppose the five blocks $T_1,T_2,T_3,P_1,P_2$ form one shared-cut cluster.
Then $\spE(G)>|V(G)|$.
\end{proposition}

\begin{proof}
For each of the thirty-six incidence trees resolved by
Lemma~\ref{lem:census}, split the certified cycle into proper consecutive
intervals, assign one interval to each incidence-tree branch, and assign every
off-core tree by its unique attachment.  The resulting connected induced
territories have exactly the branch multisets in the corresponding packet
table, whose total surplus is positive.  Equation \eqref{eq:super} settles all
thirty-six.

In exceptions 1--3, both pentagons are leaf cycle nodes and the three
triangles remain a connected shared-cut cluster after the pentagons are
opened.  Choose a private vertex $v_i$ on each $P_i$.  Let $F_i$ contain
$v_i$ and every off-cycle tree branch rooted there, and put all remaining
vertices in $H$.  Each $F_i$ is a tree, so $\sig(F_i)=-1$.  Each path
$P_i-v_i$ contains the unique shared cut of its leaf pentagon, so $H$ is
connected.  Its only cycles are the three triangles, which form one connected
shared-cut cluster.  At least two meet, the packing number is at most two, and
\eqref{eq:packet-ttt} gives $\sig(H)>2$.  Therefore
\[
 \sig(G)\geq\sig(H)+\sig(F_1)+\sig(F_2)>2-1-1=0.
\]

In exception 4, cyclically order the four marked vertices on the hub $P_1$.
The mark belonging to $P_2$ has a cyclic neighbor belonging to some triangle,
say $T_1$.  Split $P_1$ into three proper consecutive intervals: one contains
those two neighboring marks, and each of the other two contains the mark of
$T_2$ or $T_3$.  The induced territories have types $TP$, $T$, and $T$.
Consequently
\[
 \sig(G)>(1-\delta)+0+0>0.
\]
All forty colored incidence trees are covered.
\end{proof}

\section{Main theorem}

\begin{theorem}\label{thm:main}
Every connected pentacyclic cactus $G$ on $n$ vertices satisfies
\[
 \boxed{\spE(G)>n}.
\]
\end{theorem}

\begin{proof}
The exact DNN estimate \eqref{eq:dnn} proves the assertion unless the cycle
multiset is $TTTTQ$ or $TTTPP$.  If the shared-cut graph is disconnected,
Propositions~\ref{prop:disconnected-ttttq} and
\ref{prop:disconnected-tttpp} settle the two residual families, using the
complete eleven-row and fifteen-row colored partition audits.  If it is
connected, Propositions~\ref{prop:shared-ttttq} and
\ref{prop:shared-tttpp} apply.  These alternatives exhaust all cycle lengths
and all incidences.  Lemma~\ref{lem:territories} and every interval split
explicitly allow arbitrary connector trees and arbitrary attached trees.
Every spectral comparison is applied to vertex-disjoint induced subgraphs;
neither edge monotonicity nor tree relocation is used.
\end{proof}

Since a connected pentacyclic graph has $|E(G)|=n+4$, the theorem proves the
strict conclusion of the Akbari--Kumar--Mohar--Pragada--Zhang conjecture
\cite[Conjecture~1.2]{AKMPZ} for every pentacyclic cactus.

\appendix
\section{Exact reproduction and local proof objects}

From the repository root, run
\begin{verbatim}
python research/pentacyclic-tttpp-incidence-census.py
\end{verbatim}
The expected output is
\begin{verbatim}
colored trees by cut count: {1: 1, 2: 7, 3: 18, 4: 14}
one-cycle split exceptions:
  c=1: ((0, 5), (1, 5), (2, 5), (3, 5), (4, 5))
  c=2: ((0, 5), (0, 6), (1, 5), (2, 5), (3, 5), (4, 6))
  c=3: ((0, 5), (0, 6), (1, 5), (1, 7), (2, 5), (3, 6), (4, 7))
  c=4: ((0, 5), (1, 6), (2, 7), (3, 5), (3, 6), (3, 7),
        (3, 8), (4, 8))
\end{verbatim}
The script contains internal assertions for these counts, for the resolved
counts $\{2:6,3:17,4:13\}$, and for the exact four exception tuples.  The
local manuscripts supplying the stated dependencies are
\begin{verbatim}
sharp-cactus-dnn/paper.tex
all-bicyclic-cacti/paper.tex
all-tricyclic-cacti/paper.tex
all-tetracyclic-cacti/paper.tex
packing-two-square-energy/paper.tex
\end{verbatim}
The audited source proof objects integrated into the present manuscript are
\begin{verbatim}
research/pentacyclic-cactus-complete-audit-2026-07-26.md
research/pentacyclic-fully-shared-3333q-2026-07-26.md
research/pentacyclic-tttpp-fully-shared-2026-07-26.md
research/pentacyclic-tttpp-incidence-census.py
\end{verbatim}

\begin{thebibliography}{99}
\bibitem{AKMP}
S.~Akbari, H.~Kumar, B.~Mohar, and S.~Pragada,
\emph{A linear lower bound for the square energy of graphs},
Electron. J. Combin. 32(3) (2025), Paper No. P3.53.

\bibitem{AKMPZ}
S.~Akbari, H.~Kumar, B.~Mohar, S.~Pragada, and S.~Zhang,
\emph{Refinement of a conjecture on positive square energy of graphs},
arXiv:2506.07264, 2025.

\bibitem{Bicyclic}
Companion manuscript, \emph{Positive Square Energy of Every Bicyclic Cactus},
2026; repository path \texttt{all-bicyclic-cacti/paper.tex}.

\bibitem{Tricyclic}
Companion manuscript, \emph{Positive Square Energy of Every Tricyclic Cactus},
2026; repository path \texttt{all-tricyclic-cacti/paper.tex}.

\bibitem{Tetracyclic}
Companion manuscript, \emph{Positive Square Energy of Every Tetracyclic
Cactus}, 2026; repository path \texttt{all-tetracyclic-cacti/paper.tex}.

\bibitem{DNN}
Companion manuscript, \emph{The Optimized Liu--Tang--Zhang Constant for Cactus
Graphs}, 2026; repository path \texttt{sharp-cactus-dnn/paper.tex}.

\bibitem{PackingTwo}
Companion manuscript, \emph{Square Energy from Cycle Packing and Block-Graph
Territories}, 2026; repository path
\texttt{packing-two-square-energy/paper.tex}.

\bibitem{LTZ}
Y.~Liu, Q.~Tang, and S.~Zhang,
\emph{The positive and negative square-energy conjecture},
arXiv:2607.18031, 2026.
\end{thebibliography}
\end{document}
